\documentclass{article}
\usepackage{arxiv}
\usepackage{amsmath,amssymb,amsfonts}
\usepackage{natbib}
\usepackage{booktabs}
\usepackage{graphicx}
\usepackage{hyperref}
\usepackage{url}
\usepackage{xcolor}

\title{Confidence-Based Model Routing via Zoom Consistency for GUI Grounding}

\author{
  Keon Kim \\
  Om Labs \\
  \texttt{keon@omlabs.xyz} \\
  \And
  Krish Chelikavada \\
  Om Labs \\
  \texttt{krish@omlabs.xyz} \\
}

\begin{document}

\maketitle

\begin{abstract}
GUI grounding, the task of predicting click coordinates from screenshots and natural language instructions, is a critical capability for computer-use agents. Current state-of-the-art approaches use specialized vision-language models (VLMs) with multi-step zoom-in refinement pipelines. We observe that the zoom refinement step contains a free confidence signal: when a model's step-2 prediction lands near the center of the zoomed crop, the model's initial localization was already accurate. We call this signal \emph{zoom consistency} and show it correlates monotonically with prediction correctness. We use this signal to route between a specialized GUI grounding model and a general-purpose VLM, selecting the more confident prediction per sample without any training. On ScreenSpot-Pro, our method achieves \textbf{80.9\%} accuracy, a new state of the art. Code is available at \url{https://github.com/omxyz/zoom-consistency-routing}.
\end{abstract}

\section{Introduction}

GUI grounding is the task of identifying the pixel coordinates of a UI element described by a natural language instruction in a screenshot. It is a foundational capability for autonomous computer-use agents. Recent benchmarks like ScreenSpot-Pro~\cite{li2025screenspotpro} evaluate this capability on professional software (CAD tools, IDEs, video editors) with high-resolution screenshots.

The current best approach on ScreenSpot-Pro uses KV-Ground-8B~\cite{kvground}, a specialized VLM combined with a 2-step zoom-in strategy: predict a rough location on the full screenshot, crop and resize around that prediction, then re-predict on the zoomed view for refinement. This achieves 80.5\% accuracy.

We make three observations:

\textbf{First}, the zoom refinement step contains a natural confidence signal. When the model's step-2 prediction is close to the center of the zoomed crop, step-1 was already accurate and the zoom merely confirmed it. When step-2 points far from the crop center, step-1 was off, and step-2 is attempting a large correction that often fails. We formalize this as \emph{zoom consistency}.

\textbf{Second}, different VLMs have complementary failure modes. A specialized model like KV-Ground-8B excels on professional engineering tools, while a general-purpose model like Qwen3.5-27B~\cite{qwen35} occasionally succeeds on consumer interfaces where the specialist fails. An oracle that always picks the correct model achieves 85.1\%, revealing 5 percentage points of untapped potential.

\textbf{Third}, zoom consistency enables training-free routing between models. By running both models independently and selecting the one with lower zoom consistency per sample, we capture a portion of this complementary potential without any additional training or labeled routing data.

Our contributions:
\begin{enumerate}
    \item We identify \emph{zoom consistency} as a free, per-sample confidence signal inherent to multi-step zoom-in pipelines for GUI grounding.
    \item We validate that zoom consistency correlates monotonically with prediction accuracy across 1,581 professional GUI samples.
    \item We propose a training-free heterogeneous ensemble that routes between a specialized and a general-purpose VLM using zoom consistency, achieving 80.9\% on ScreenSpot-Pro.
\end{enumerate}

\section{Related Work}

\paragraph{GUI Grounding.}
GUI grounding has evolved from rule-based approaches to VLM-based methods. SeeClick~\cite{seeclick} established the first realistic GUI grounding benchmark. OS-Atlas~\cite{osatlas} and UGround~\cite{uground} provided large-scale training data across platforms. GUI-Actor~\cite{guiactor} introduced coordinate-free grounding using set-of-mark prompting. UI-TARS~\cite{uitars} scaled to 72B parameters. MEGA-GUI~\cite{megagui} proposed multi-stage coarse-to-fine grounding. KV-Ground-8B~\cite{kvground} achieved state-of-the-art results through multi-stage fine-tuning: Qwen3-VL-8B was first trained as GUI-Owl-1.5~\cite{guiowl,mplugowl} on millions of GUI agent samples, then further specialized for coordinate grounding.

\paragraph{Zoom-In Strategies.}
Multi-step zoom-in is a widely adopted strategy for high-resolution GUI grounding. ZoomClick~\cite{zoomclick} systematically evaluated zoom refinement on ScreenSpot-Pro. LASER~\cite{laser} introduced preference optimization for zoom region selection. ScreenSeekeR~\cite{screenseeker} used cascaded visual search with zoom. AdaZoom-GUI~\cite{adazoomgui} introduced adaptive zoom with instruction refinement. GUI-G$^2$~\cite{guig2} proposed Gaussian reward modeling for zoom-based grounding RL. We do not propose a new zoom strategy. We extract a confidence signal from existing zoom pipelines.

\paragraph{Confidence Estimation and Model Routing.}
GUI-RC~\cite{guirc} introduced region consistency as a test-time confidence signal, using spatial voting across multiple predictions. Adaptive VLM Routing~\cite{avr} proposed routing between VLMs based on action difficulty estimation. Training-Free Uncertainty Guidance~\cite{uncertaintyguidance} uses model uncertainty to select visual inputs. Adaptive Chain-of-Focus~\cite{adaptivecof} dynamically decides when to zoom. Our approach requires no additional sampling, no difficulty estimation, and no training. The confidence signal is a byproduct of the zoom pipeline that models already run.

\section{Method}

\subsection{Background: 2-Step Zoom-In Pipeline}

Given a screenshot $I$ of size $W \times H$ and a natural language instruction $q$, the standard zoom-in pipeline works as follows:

\textbf{Step 1.} Run the VLM on the full image to obtain a rough prediction $\hat{p}_1 = (x_1, y_1)$ in a normalized $1000\times1000$ coordinate space.

\textbf{Step 2.} Compute a crop box centered on $\hat{p}_1$ with crop ratio $r$ (typically 0.5), crop the image, resize the crop back to the original resolution, and re-run the VLM to obtain $\hat{p}_2 = (x_2, y_2)$ in the crop's $1000\times1000$ space.

\textbf{Step 3.} Remap $\hat{p}_2$ from crop coordinates back to full-image coordinates to produce the final prediction.

\subsection{Zoom Consistency}

We define \emph{zoom consistency} as the Euclidean distance between the step-2 prediction and the center of the crop in the model's coordinate space:
\begin{equation}
    c = \sqrt{(x_2 - 500)^2 + (y_2 - 500)^2}
\end{equation}
where $(x_2, y_2)$ is the step-2 prediction in the $1000\times1000$ crop space, and $(500, 500)$ is the crop center.

\textbf{Intuition.} After cropping around $\hat{p}_1$, the predicted location should be near the center of the crop if $\hat{p}_1$ was accurate. A low $c$ (close to center) indicates step-1 was already correct, and step-2 merely confirms it. A high $c$ (far from center) indicates step-1 was off, and step-2 is attempting a large spatial correction.

\textbf{Key property:} Zoom consistency is \emph{free}: it requires no additional model calls, no sampling, no temperature variation. It is computed from values already produced by the standard pipeline.

\subsection{Consistency-Based Routing}

Given two models $M_A$ (specialist) and $M_B$ (generalist), the routing procedure is:

\begin{enumerate}
    \item Run $M_A$'s full 2-step zoom pipeline $\rightarrow$ prediction $\hat{p}_A$, consistency $c_A$
    \item Run $M_B$'s full 2-step zoom pipeline $\rightarrow$ prediction $\hat{p}_B$, consistency $c_B$
    \item Select the final prediction:
\end{enumerate}
\begin{equation}
    \hat{p} = \begin{cases} \hat{p}_A & \text{if } c_A \leq c_B \\ \hat{p}_B & \text{otherwise} \end{cases}
\end{equation}

No training is required. No threshold tuning is needed. The router picks whichever model is more self-consistent in its zoom refinement.

\subsection{Model Selection}

We pair two models with complementary strengths:

\textbf{KV-Ground-8B} (specialist): A Qwen3-VL-8B model fine-tuned through GUI-Owl-1.5 on millions of GUI screenshots, then further specialized for coordinate grounding. Strong on professional tools (CAD, IDEs, creative software).

\textbf{Qwen3.5-27B} (generalist): A recent unified multimodal model with strong general visual understanding. It is not fine-tuned for GUI grounding specifically, but occasionally succeeds where KV-Ground fails, particularly on consumer/OS interfaces.

\section{Experimental Setup}

\paragraph{Benchmark.}
We evaluate on \textbf{ScreenSpot-Pro}~\cite{li2025screenspotpro,screenspotpro_repo}, the most challenging GUI grounding benchmark. It contains 1,581 samples across 26 professional applications, 6 industry categories (CAD, Creative, Dev, Office, OS, Scientific), and 3 operating systems. Each sample consists of a high-resolution screenshot, a natural language instruction, and a bounding box annotation. Accuracy is measured as point-in-box: whether the predicted coordinate falls within the annotated bounding box.

\paragraph{Implementation Details.}
KV-Ground-8B uses \texttt{vocaela/KV-Ground-8B-BaseGuiOwl1.5-0315} loaded in bf16 with no resolution cap (\texttt{max\_pixels}$=99{,}999{,}999$). Qwen3.5-27B uses \texttt{cyankiwi/Qwen3.5-27B-AWQ-4bit} loaded via compressed\_tensors. Both use SDPA attention. The zoom pipeline uses 2 steps with \texttt{crop\_ratio}$=0.5$ and greedy decoding (temperature$=0$). All experiments run on NVIDIA H200 141GB.

\paragraph{Baselines.}
We compare against: KV-Ground-8B (single model, 2-step zoom), the current leaderboard \#1; Qwen3.5-27B (single model, 2-step zoom); stage split (KV-Ground step-1, Qwen step-2); midpoint fusion; vote with agreement; KV-default fallback; and oracle (always picks the correct model).

\section{Results}

\subsection{Main Results}

\begin{table}[h]
\centering
\caption{Comparison of routing strategies on ScreenSpot-Pro (1,581 samples).}
\label{tab:main}
\begin{tabular}{lccr}
\toprule
\textbf{Method} & \textbf{Accuracy} & \textbf{Correct} & \textbf{vs.\ KV} \\
\midrule
\textbf{Consistency Router (ours)} & \textbf{80.9\%} & \textbf{1,279} & \textbf{+0.8\%} \\
KV-Ground-8B (baseline) & 80.1\% & 1,266 & --- \\
Stage split KV$\rightarrow$Qwen & 78.7\% & 1,245 & $-$1.3\% \\
Vote agree $T{=}50$ & 76.9\% & 1,215 & $-$3.2\% \\
Midpoint fusion & 75.7\% & 1,196 & $-$4.4\% \\
Step-1 vote centroid & 72.5\% & 1,146 & $-$7.6\% \\
KV fallback $T{=}300$ & 69.6\% & 1,100 & $-$10.5\% \\
Qwen3.5-27B only & 60.9\% & 962 & $-$19.2\% \\
\midrule
Oracle (upper bound) & 85.1\% & 1,345 & +5.0\% \\
\bottomrule
\end{tabular}
\end{table}

Our consistency router achieves 80.9\%, a new state of the art on ScreenSpot-Pro. All other ensemble strategies performed \emph{worse} than the KV-Ground baseline, showing that naive ensembling hurts in this domain and that confidence-based routing is essential.

\subsection{Per-Category Results}

\begin{table}[h]
\centering
\caption{Per-category results of the zoom consistency router.}
\label{tab:category}
\begin{tabular}{lrrrr}
\toprule
\textbf{Category} & \textbf{$n$} & \textbf{Icon} & \textbf{Text} & \textbf{Avg} \\
\midrule
Office & 230 & 75.5\% & 94.9\% & 90.4\% \\
Scientific & 254 & 70.9\% & 93.8\% & 83.9\% \\
Dev & 299 & 73.1\% & 93.5\% & 83.6\% \\
CAD & 261 & 51.6\% & 87.8\% & 78.9\% \\
Creative & 341 & 62.2\% & 85.4\% & 75.7\% \\
OS & 196 & 56.2\% & 87.8\% & 73.5\% \\
\midrule
\textbf{Overall} & \textbf{1,581} & \textbf{65.6\%} & \textbf{90.4\%} & \textbf{80.9\%} \\
\bottomrule
\end{tabular}
\end{table}

The router's gains are distributed across categories, with no single category driving the improvement.

\subsection{Routing Statistics and Per-Application Analysis}

The router selects KV-Ground for 64.9\% of samples and Qwen3.5-27B for 35.1\%. The routing distribution varies by category (Table~\ref{tab:routing}).

\begin{table}[h]
\centering
\caption{Router model selection by category. Qwen is selected more often for general-purpose software and less for specialist professional tools.}
\label{tab:routing}
\begin{tabular}{lrrr}
\toprule
\textbf{Category} & \textbf{$n$} & \textbf{KV selected} & \textbf{Qwen selected} \\
\midrule
CAD & 261 & 86.6\% & 13.4\% \\
Dev & 299 & 64.9\% & 35.1\% \\
Scientific & 254 & 63.8\% & 36.2\% \\
Creative & 341 & 61.9\% & 38.1\% \\
OS & 196 & 57.1\% & 42.9\% \\
Office & 230 & 52.6\% & 47.4\% \\
\bottomrule
\end{tabular}
\end{table}

The router's behavior reflects domain-specific confidence patterns. For CAD applications (AutoCAD, Inventor, SolidWorks), KV-Ground is selected 86.6\% of the time. KV-Ground was trained on millions of professional GUI screenshots, so its zoom consistency is high on these specialist interfaces. For Office applications (Word, Excel, PowerPoint), Qwen is selected 47.4\% of the time, nearly equal to KV-Ground. These consumer interfaces are well-represented in Qwen3.5's general pretraining data.

Table~\ref{tab:perapps} shows the per-application breakdown for the 15 applications where the router changes the outcome relative to KV-Ground.

\begin{table}[h]
\centering
\caption{Per-application accuracy where the router differs from KV-Ground baseline.}
\label{tab:perapps}
\small
\begin{tabular}{lrrrrr}
\toprule
\textbf{Application} & \textbf{$n$} & \textbf{KV} & \textbf{Router} & \textbf{$\Delta$} & \textbf{Qwen \%} \\
\midrule
photoshop & 51 & 76.5\% & 84.3\% & +7.8\% & 49.0\% \\
macos\_common & 65 & 75.4\% & 80.0\% & +4.6\% & 46.2\% \\
illustrator & 31 & 83.9\% & 87.1\% & +3.2\% & 25.8\% \\
blender & 71 & 76.1\% & 78.9\% & +2.8\% & 40.8\% \\
pycharm & 78 & 78.2\% & 80.8\% & +2.6\% & 43.6\% \\
windows\_common & 81 & 65.4\% & 67.9\% & +2.5\% & 39.5\% \\
vmware & 41 & 92.7\% & 95.1\% & +2.4\% & 58.5\% \\
stata & 49 & 77.6\% & 79.6\% & +2.0\% & 36.7\% \\
android\_studio & 80 & 78.7\% & 80.0\% & +1.3\% & 17.5\% \\
word & 84 & 92.9\% & 94.0\% & +1.2\% & 50.0\% \\
matlab & 93 & 90.3\% & 91.4\% & +1.1\% & 54.8\% \\
\midrule
vivado & 80 & 77.5\% & 76.2\% & $-$1.3\% & 25.0\% \\
quartus & 45 & 88.9\% & 86.7\% & $-$2.2\% & 28.9\% \\
vscode & 55 & 85.5\% & 81.8\% & $-$3.6\% & 36.4\% \\
davinci & 44 & 79.5\% & 75.0\% & $-$4.5\% & 22.7\% \\
\bottomrule
\end{tabular}
\end{table}

The router improves accuracy on 11 applications and reduces it on 4. The largest gain is on Photoshop (+7.8\%), where Qwen3.5-27B's general visual understanding helps with Photoshop's standard toolbar and menu layouts. The largest regression is on DaVinci Resolve ($-$4.5\%), where the router occasionally selects Qwen on samples that KV-Ground handles correctly. The net effect across all 26 applications is positive (+0.8\%).

\subsection{Icon vs.\ Text Analysis}

\begin{table}[h]
\centering
\caption{Router gains on icon vs.\ text UI elements.}
\label{tab:uitype}
\begin{tabular}{lrrrr}
\toprule
\textbf{UI Type} & \textbf{$n$} & \textbf{KV baseline} & \textbf{Router} & \textbf{$\Delta$} \\
\midrule
Icon & 604 & 64.6\% & 65.6\% & +1.0\% \\
Text & 977 & 89.7\% & 90.4\% & +0.7\% \\
\bottomrule
\end{tabular}
\end{table}

Gains are balanced between icon and text targets. The router does not disproportionately benefit one UI element type over the other.

\section{Analysis}

\subsection{Zoom Consistency as a Confidence Signal}

We bucket samples by zoom consistency and measure accuracy for each model independently (Table~\ref{tab:consistency}).

\begin{table}[h]
\centering
\caption{KV-Ground accuracy by zoom consistency bucket. Lower consistency (closer to crop center) correlates with higher accuracy.}
\label{tab:consistency}
\begin{tabular}{lrr}
\toprule
\textbf{Consistency $c$} & \textbf{$n$} & \textbf{Accuracy} \\
\midrule
$< 30$ (very confident) & 464 & 87.1\% \\
$30$--$80$ & 101 & 85.1\% \\
$80$--$150$ & 106 & 76.4\% \\
$150$--$250$ & 173 & 79.8\% \\
$\geq 250$ (uncertain) & 737 & 75.6\% \\
\bottomrule
\end{tabular}
\end{table}

The correlation is monotonic: samples where step-2 barely moved from center (consistency $< 30$) have 87.1\% accuracy, while samples where step-2 made large corrections ($\geq 250$) have 75.6\%. This 11.5-point gap confirms that zoom consistency is a meaningful confidence signal.

Qwen3.5-27B shows a similar pattern: 90.2\% accuracy at consistency $< 30$ versus 80.7\% at consistency $\geq 250$, a 9.5-point gap. The weaker gap for Qwen is expected because Qwen was not trained on the zoom pipeline and its step-2 behavior is less calibrated.

\subsection{Why the Router is Adaptive}

The router adapts its behavior based on domain-specific confidence patterns. On CAD applications, KV-Ground's zoom consistency is typically low (the specialist is confident on specialist tools), so the router selects KV-Ground 86.6\% of the time. On Office and OS applications, Qwen3.5-27B's zoom consistency is often competitive with or lower than KV-Ground's, reflecting Qwen's familiarity with consumer interfaces from pretraining. The router responds by selecting Qwen 42--47\% of the time on these categories.

This adaptive behavior emerges from the consistency signal alone, with no application labels or category metadata provided to the router.

\subsection{Why Other Ensemble Strategies Fail}

\textbf{Stage split} (KV step-1, Qwen step-2) forces Qwen to refine every sample's zoomed crop, even when KV-Ground would have been better at refinement. On ScreenSpot-Pro, KV-Ground's GUI-specific training makes it a better refiner on 78\% of samples.

\textbf{Midpoint fusion} averages two coordinate predictions, often pulling a correct prediction out of the bounding box. GUI bounding boxes are small (median area 0.1\% of the image), so even a small perturbation from averaging causes a miss.

\textbf{Vote with agreement} provides no mechanism to exploit disagreement constructively. When models disagree, defaulting to KV-Ground misses Qwen's unique wins.

\textbf{KV fallback with threshold} routes to Qwen when KV is uncertain. This fails because Qwen's overall accuracy (60.9\%) is far below KV's (80.1\%). Even when KV is uncertain, it is still more likely correct than Qwen.

The consistency router succeeds because it can always choose KV-Ground when KV is more confident, uses a per-sample signal rather than static rules, and measures each model's self-consistency independently through its own zoom pipeline.

\subsection{Oracle Analysis}

Of the 1,581 samples, the oracle (which always picks the correct model) achieves 85.1\% (1,345 correct). The breakdown:

\begin{itemize}
    \item \textbf{Both correct}: 962 samples (60.8\%)
    \item \textbf{Only KV correct}: 304 samples (19.2\%)
    \item \textbf{Only Qwen correct}: 79 samples (5.0\%)
    \item \textbf{Both wrong}: 236 samples (14.9\%)
\end{itemize}

The router's task is to capture as many of the 79 Qwen-unique wins as possible without losing the 304 KV-unique wins. It captures 13 of the 79 (16.5\% capture rate) while losing minimal KV-unique wins, yielding a net gain of +13 correct samples.

\subsection{Screening vs.\ Full Evaluation}

We initially validated on a 200-sample screening subset before running the full 1,581-sample evaluation (Table~\ref{tab:screening}).

\begin{table}[h]
\centering
\caption{Router gain on screening (200 samples) vs.\ full evaluation (1,581 samples).}
\label{tab:screening}
\begin{tabular}{lrr}
\toprule
\textbf{Metric} & \textbf{Screening} & \textbf{Full} \\
\midrule
KV-Ground baseline & 84.5\% & 80.1\% \\
Consistency router & 86.0\% & 80.9\% \\
Router gain & +1.5\% & +0.8\% \\
Oracle & 89.5\% & 85.1\% \\
\bottomrule
\end{tabular}
\end{table}

The router gain shrinks from +1.5\% on screening to +0.8\% on full evaluation. The full dataset includes harder professional tools where both models struggle and the oracle advantage is smaller. Screening results can overstate improvements by roughly $2\times$.

\section{Limitations}

\begin{enumerate}
    \item \textbf{Computational cost.} The router runs two full models (2 forward passes each, 4 total), doubling inference time compared to a single model.
    \item \textbf{Modest gains.} The improvement (+0.8\% on full eval) is real but small. A stronger generalist model could increase the capture rate beyond 16.5\% of the oracle advantage.
    \item \textbf{Generalist model quality.} Qwen3.5-27B-AWQ achieves 60.9\% standalone accuracy, limiting the oracle ceiling. A GUI-trained generalist could significantly increase the complementary advantage.
    \item \textbf{Zoom consistency limitations.} The signal is strongest for KV-Ground (11.5-point gap) and weaker for Qwen (9.5-point gap). Models trained specifically for zoom pipelines may have artificially calibrated consistency, reducing discriminative power.
    \item \textbf{Per-application regressions.} The router reduces accuracy on 4 of 26 applications (DaVinci $-$4.5\%, VSCode $-$3.6\%, Quartus $-$2.2\%, Vivado $-$1.3\%). These regressions are outweighed by gains elsewhere, but a production deployment may need per-application safeguards.
\end{enumerate}

\section{Conclusion}

We introduced \emph{zoom consistency}, the distance between a model's zoom-refinement prediction and the crop center, as a free, per-sample confidence signal for GUI grounding. By routing between KV-Ground-8B and Qwen3.5-27B based on which has lower zoom consistency, we achieve 80.9\% on ScreenSpot-Pro without any additional training. The router adapts its behavior by domain: selecting the specialist 87\% of the time on CAD applications and the generalist up to 47\% of the time on consumer software, all from the consistency signal alone. The zoom consistency signal is general and could be applied to any multi-step refinement pipeline.

\bibliography{references}
\bibliographystyle{plainnat}

\end{document}
